\documentclass{amsart}
\usepackage{amsmath,amssymb,amsthm,hyperref}
\newtheorem{theorem}{Theorem}
\newtheorem{lemma}{Lemma}
\newtheorem{proposition}{Proposition}
\theoremstyle{remark}
\newtheorem{remark}{Remark}
\begin{document}

\title{The functional equation is equivalent to the Collatz conjecture}
\maketitle

\section*{Statement}
Let $\mathbb{D} = \{z \in \mathbb{C} : |z| < 1\}$ and $\lambda = e^{2\pi i/3}$.
We study the analytic functions $\varphi:\mathbb{D}\to\mathbb{C}$ satisfying
\begin{equation}\label{eq:FE}
3z\,\varphi(z^3) \;=\; 3z\,\varphi(z^6) \;+\; \varphi(z^2) \;+\; \lambda\,\varphi(\lambda z^2) \;+\; \lambda^2 \varphi(\lambda^2 z^2),
\qquad z\in\mathbb{D}.
\end{equation}

The question asks whether \emph{every} analytic solution has the form
$\varphi(z)=a+\dfrac{bz}{1-z}$, $a,b\in\mathbb{C}$.

\medskip
\noindent\textbf{Main result.} \emph{The proposed statement is logically equivalent to the
Collatz ($3n+1$) conjecture.} Consequently, with current knowledge it can neither be proved
nor disproved; below we give a complete and rigorous reduction. (We also record precisely
which implication needs which form of Collatz.)

\section{Reduction of the functional equation to a coefficient recursion}

Since $\varphi$ is analytic on $\mathbb{D}$ it has a power series
\begin{equation}\label{eq:series}
\varphi(z)=\sum_{n\ge 0}c_n z^n,\qquad \limsup_{n\to\infty}|c_n|^{1/n}\le 1,
\end{equation}
converging on $\mathbb{D}$. All the substitutions $z\mapsto z^2,z^3,z^6,\lambda z^2,\lambda^2 z^2$
keep $z$ inside $\mathbb{D}$, so each side of \eqref{eq:FE} is an absolutely convergent power
series on $\mathbb{D}$, and \eqref{eq:FE} holds if and only if the coefficients of $z^k$ agree
for every $k\ge 0$.

\begin{lemma}\label{lem:sym}
For every $n\ge 0$,
\[
1+\lambda^{\,n+1}+\lambda^{\,2(n+1)}=
\begin{cases}
3,& n\equiv 2 \pmod 3,\\[2pt]
0,& \text{otherwise.}
\end{cases}
\]
\end{lemma}
\begin{proof}
This is the standard root-of-unity filter: $1+\omega+\omega^2=0$ for any cube root of unity
$\omega\ne 1$, while it equals $3$ when $\omega=1$. Here $\omega=\lambda^{\,n+1}$, which equals
$1$ exactly when $3\mid n+1$, i.e.\ $n\equiv 2\pmod 3$.
\end{proof}

\begin{lemma}\label{lem:terms}
With $\varphi$ as in \eqref{eq:series},
\[
\varphi(z^2)+\lambda\,\varphi(\lambda z^2)+\lambda^2\varphi(\lambda^2 z^2)
=3\!\!\sum_{\substack{m\ge 0\\ m\equiv 2\,(3)}}\!\! c_m\,z^{2m}.
\]
\end{lemma}
\begin{proof}
The left side equals
$\sum_{m\ge 0} c_m\bigl(z^{2m}+\lambda(\lambda z^2)^m+\lambda^2(\lambda^2 z^2)^m\bigr)
=\sum_{m\ge0}c_m z^{2m}\bigl(1+\lambda^{m+1}+\lambda^{2(m+1)}\bigr)$,
and we apply Lemma~\ref{lem:sym}.
\end{proof}

Substituting Lemma~\ref{lem:terms} and the expansions
$3z\,\varphi(z^3)=3\sum_{n\ge0}c_n z^{3n+1}$ and
$3z\,\varphi(z^6)=3\sum_{n\ge0}c_n z^{6n+1}$ into \eqref{eq:FE} and dividing by $3$, the
functional equation \eqref{eq:FE} is equivalent to the formal/identity statement
\begin{equation}\label{eq:three}
\sum_{n\ge0}c_n z^{3n+1}
=\sum_{n\ge0}c_n z^{6n+1}
+\!\!\sum_{\substack{m\ge0\\ m\equiv2\,(3)}}\!\! c_m\,z^{2m}.
\end{equation}

\begin{lemma}\label{lem:exponents}
Every monomial appearing on either side of \eqref{eq:three} has exponent $\equiv 1\pmod 3$.
Moreover, for each $n\ge 0$ the coefficient of $z^{3n+1}$ in \eqref{eq:three} yields exactly the
single relation
\begin{equation}\label{eq:recursion}
c_n \;=\;
\begin{cases}
c_{\,n/2}, & n\ \text{even},\\[2pt]
c_{\,(3n+1)/2}, & n\ \text{odd},
\end{cases}
\end{equation}
and these are \emph{all} the equations imposed by \eqref{eq:FE}.
\end{lemma}
\begin{proof}
On the left the exponents are $3n+1\equiv1\pmod3$. The first sum on the right has exponents
$6n+1\equiv1\pmod3$. In the last sum the exponents are $2m$ with $m\equiv2\pmod3$, hence
$2m\equiv 4\equiv1\pmod 3$. Thus every exponent is $\equiv1\pmod3$, and matching the coefficient
of $z^{k}$ is nontrivial only for $k=3n+1$.

Fix $n\ge0$ and put $k=3n+1$. The left side of \eqref{eq:three} contributes the single term
$c_n$ (from its index $n$, as $3n+1$ is injective in $n$). On the right:

\emph{First sum:} it contributes $c_j$ when $6j+1=3n+1$, i.e.\ $6j=3n$, i.e.\ $n=2j$. This
happens iff $n$ is even, in which case $j=n/2$ and the contribution is $c_{n/2}$. (Each value
$6j+1$ occurs for a unique $j$.)

\emph{Last sum:} it contributes $c_m$ when $2m=3n+1$ with $m\equiv2\pmod3$. The equation
$2m=3n+1$ has an integer solution $m$ iff $3n+1$ is even, i.e.\ iff $n$ is odd; then
$m=(3n+1)/2$. Writing $n=2t+1$ gives $m=3t+2\equiv2\pmod3$, so the congruence condition is
automatically satisfied. Hence for odd $n$ the contribution is exactly $c_{(3n+1)/2}$.

Finally, for even $n$ we have $3n+1$ odd, so the last sum does not contribute; for odd $n$ we
have $3n+1\equiv 4\pmod 6$, so the first sum does not contribute. Therefore the coefficient
identity for $z^{3n+1}$ reads precisely \eqref{eq:recursion}. As $n$ ranges over all
nonnegative integers, $k=3n+1$ ranges over all admissible exponents, so \eqref{eq:recursion}
for all $n\ge0$ is the complete content of \eqref{eq:FE}.
\end{proof}

\begin{remark}
For $n=0$ relation \eqref{eq:recursion} reads $c_0=c_0$, an empty constraint; so $c_0$ is a free
parameter (it will play the role of $a$).
\end{remark}

\section{The Collatz (Syracuse) map appears}

Define the \emph{Syracuse map} $T:\mathbb{Z}_{\ge1}\to\mathbb{Z}_{\ge1}$ by
\[
T(n)=
\begin{cases}
n/2, & n\ \text{even},\\
(3n+1)/2, & n\ \text{odd}.
\end{cases}
\]
(Note $T$ maps positive integers to positive integers, and $T(1)=2$, $T(2)=1$.) Lemma
\ref{lem:exponents} says exactly that an analytic $\varphi=\sum c_n z^n$ solves \eqref{eq:FE}
if and only if
\begin{equation}\label{eq:orbit}
c_n=c_{T(n)}\qquad\text{for every } n\ge1,
\end{equation}
with $c_0\in\mathbb{C}$ arbitrary.

Thus \emph{the solution space is the space of functions $n\mapsto c_n$ on $\mathbb{Z}_{\ge1}$
that are constant on the forward orbits of $T$, together with a free constant $c_0$, subject only
to the growth condition} $\limsup|c_n|^{1/n}\le1$.

Introduce the equivalence relation $\sim$ on $\mathbb{Z}_{\ge1}$ generated by $n\sim T(n)$; its
classes are the weakly connected components of the functional graph of $T$. Equation
\eqref{eq:orbit} holds for an analytic $\varphi$ \emph{iff} the sequence $(c_n)_{n\ge1}$ is
constant on each $\sim$-class (constancy on a class is equivalent to constancy along every edge
$n\to T(n)$, since the classes are generated by those edges).

\begin{proposition}\label{prop:correspondence}
There is a bijection between analytic solutions of \eqref{eq:FE} and pairs
$\bigl(c_0,\ (\gamma_C)_{C}\bigr)$, where $c_0\in\mathbb{C}$ is arbitrary and
$(\gamma_C)_C$ is a bounded\footnote{Boundedness suffices for analyticity on $\mathbb{D}$ because
$\sup_n|c_n|<\infty$ forces $\limsup|c_n|^{1/n}\le1$; conversely the growth condition restricts
which assignments are admissible, but every \emph{bounded} assignment is admissible.}
family of complex numbers indexed by the $\sim$-classes $C$ of $\mathbb{Z}_{\ge1}$, via
$c_n=\gamma_{C(n)}$ for $n\ge1$, where $C(n)$ is the class of $n$.
\end{proposition}
\begin{proof}
Given a solution, set $c_0$ and let $\gamma_C$ be the common value of $c_n$ on each class $C$
(well defined by \eqref{eq:orbit}); boundedness of the bounded part follows from
$\limsup|c_n|^{1/n}\le1$ only in the bounded case, but any bounded family yields an analytic
function, which is all we need below. Conversely, given $c_0$ and a bounded family
$(\gamma_C)$, the sequence $c_n:=\gamma_{C(n)}$ ($n\ge1$) is bounded, hence
$\varphi(z)=c_0+\sum_{n\ge1}c_n z^n$ converges on $\mathbb{D}$ and satisfies \eqref{eq:orbit},
i.e.\ \eqref{eq:FE}. The two constructions are mutually inverse.
\end{proof}

\section{Equivalence with the Collatz conjecture}

Recall the \textbf{Collatz conjecture}: for every integer $n\ge1$ the trajectory
$n,\,T(n),\,T^2(n),\dots$ eventually reaches $1$. Equivalently, $T$ has a unique
weakly connected component on $\mathbb{Z}_{\ge1}$ (namely all integers are $\sim$ to $1$); a
counterexample is either a periodic orbit other than $\{1,2\}$ or a trajectory tending to
infinity, and in either case produces a second $\sim$-class.

\begin{theorem}\label{thm:main}
The assertion
\begin{quote}
``every analytic $\varphi:\mathbb{D}\to\mathbb{C}$ solving \eqref{eq:FE} has the form
$\varphi(z)=a+\dfrac{bz}{1-z}$''
\end{quote}
is true if and only if the Collatz conjecture is true.
\end{theorem}

\begin{proof}
First note the target family in coefficient terms: since
$\dfrac{bz}{1-z}=b\sum_{n\ge1}z^n$, the function $a+\dfrac{bz}{1-z}$ has $c_0=a$ and $c_n=b$ for
all $n\ge1$. Thus ``$\varphi$ is of the prescribed form'' is exactly the condition
\[
c_n \text{ is the same for all } n\ge1.
\tag{$\star$}
\]
By Proposition~\ref{prop:correspondence}, $(\star)$ holds for \emph{every} solution if and only
if there is only one possible value of $c_n$ across $n\ge1$ for every admissible family, which is
the case precisely when there is a \emph{single} $\sim$-class on $\mathbb{Z}_{\ge1}$.

($\Leftarrow$) Assume Collatz holds. Then every $n\ge1$ satisfies $n\sim 1$, so there is a single
class. Hence for any solution, $c_n=c_1$ for all $n\ge1$ by \eqref{eq:orbit}. Setting $a:=c_0$ and
$b:=c_1$ gives $\varphi(z)=a+b\sum_{n\ge1}z^n=a+\dfrac{bz}{1-z}$. (Convergence of the latter on
$\mathbb{D}$ is automatic, and analyticity is clear.) So the statement is true.

($\Rightarrow$) Assume Collatz fails. Then $\mathbb{Z}_{\ge1}$ has at least two distinct
$\sim$-classes; let $C_1$ be the class of $1$ (which contains $\{1,2\}$) and let $C_2\ne C_1$ be
another class. Define a bounded family by $\gamma_{C_2}=1$ and $\gamma_C=0$ for all $C\ne C_2$,
and set $c_0=0$. By Proposition~\ref{prop:correspondence} the function
\[
\varphi_\star(z)=\sum_{n:\,n\in C_2} z^n
\]
is analytic on $\mathbb{D}$ (its coefficients are $0/1$, hence bounded, so the radius of
convergence is $\ge1$) and solves \eqref{eq:FE}. But $\varphi_\star$ has $c_n=1$ for
$n\in C_2$ and $c_n=0$ for $n\in C_1$ (e.g.\ $c_1=0\ne 1=c_{m}$ for some $m\in C_2$), so
$(\star)$ fails: $\varphi_\star$ is \emph{not} of the form $a+bz/(1-z)$. Hence the statement is
false.

Combining the two implications proves the theorem.
\end{proof}

\begin{remark}[Honest status]
Because the Collatz conjecture is a famous open problem, Theorem~\ref{thm:main} shows that the
original ``prove or disprove'' question is, as of now, \emph{undecided}: a proof of the statement
would prove Collatz, and a single counterexample analytic solution (necessarily built from a
non-trivial Collatz cycle or divergent trajectory) would disprove Collatz. What \emph{is} proved
unconditionally here is:
\begin{itemize}
\item The complete classification \eqref{eq:orbit}: analytic solutions of \eqref{eq:FE} are
exactly the analytic functions whose Taylor coefficients are constant on Syracuse orbits, with
$c_0$ free.
\item The verification that $a+bz/(1-z)$ are solutions (they correspond to the constant
assignment $c_n\equiv b$, $n\ge1$, which is constant on every orbit).
\item The exact logical equivalence of the proposed uniqueness statement with the Collatz
conjecture (Theorem~\ref{thm:main}).
\end{itemize}
\end{remark}

\end{document}
